\documentclass[10pt]{article}
\usepackage[letterpaper,margin=0.9in]{geometry}
\usepackage{amsmath,amssymb,amsthm,url}
\urlstyle{same}
\urldef\repourl\url{https://github.com/KamyuRuri/AI-Math-Proofs}
\newtheorem{theorem}{Theorem}
\newcommand{\gtot}{\gamma_{\mathrm t}}
\newcommand{\ec}{\operatorname{ecc}}
\raggedbottom
\title{An 11-vertex counterexample to Graffiti.pc Conjecture 255\\
  \large Minimum order and the sharp universal constant}
\author{ChatGPT (OpenAI, GPT-6), Gemini, and JiaZhuang Wang%
  \thanks{Corresponding author. Repository:
  \repourl.}}
\date{November 27, 2026}
\begin{document}
\maketitle
\begin{abstract}
Graffiti.pc Conjecture 255 relates the total domination number of a connected
graph to the size of its center and the largest neighborhood of an edge.
Subdividing the five edges incident with one vertex of \(K_6\) gives an
11-vertex graph on which the conjecture asserts \(21\ge22\).
We prove that 11 is the minimum order of a counterexample and extend the
construction to an infinite family. We also determine the largest coefficient
that can replace the conjectured factor \(2\): it is \(1\), as shown by an
augmented family of rook graphs.
\end{abstract}

\section{The conjecture and a family of counterexamples}
Let \(G=(V,E)\) be a finite simple connected graph of order \(n\ge2\).
Write \(N(v)\) for the open neighborhood of \(v\). A total dominating set
is a set \(D\subseteq V\) for which \(\bigcup_{v\in D}N(v)=V\);
its minimum cardinality is \(\gtot(G)\). Write \(C(G)\) for the vertices
of minimum eccentricity, where \(\ec(v)=\max_{w\in V}d(v,w)\), and set
\[
 M(G)=\max_{uv\in E}|N(u)\cup N(v)|.
\]
Each endpoint of \(uv\) belongs to this union through the other endpoint.
Conjecture 255 in Ermelinda DeLaVi\~na's \emph{Written on the Wall II}
(February 23, 2007) proposes \(\gtot(G)M(G)\ge 2|C(G)|\).

\begin{theorem}
For \(t\ge2\), let \(G_t\) be obtained from \(K_{t+1}\) by subdividing
each edge at a fixed vertex once. Then
\[
 |C(G_t)|=2t+1,\qquad M(G_t)=t+2,\qquad \gtot(G_t)=3.
\]
In particular, \(2|C(G_t)|-\gtot(G_t)M(G_t)=t-4\).
\end{theorem}
\begin{proof}
Call the fixed vertex \(x\), the remaining clique \(A=\{a_1,\ldots,a_t\}\),
and the subdivision vertices \(B=\{b_1,\ldots,b_t\}\), with paths
\(x b_i a_i\). Every pair is at distance at most two. In particular,
\(b_i,b_j\) share \(x\), and for \(i\ne j\) the path from \(b_i\) to \(a_j\)
goes through \(a_i\). No vertex is universal, so all have eccentricity two.

For edges of the three types \(a_i a_j,\ a_i b_i,\ xb_i\), the neighborhood
unions are respectively
\[
 A\cup\{b_i,b_j\},\qquad A\cup\{b_i,x\},\qquad B\cup\{x,a_i\}.
\]
All have size \(t+2\). The set \(\{x,b_1,a_1\}\) totally dominates \(G_t\).
A pair that totally dominates must itself form an edge, whose neighborhood
union would then have size \(2t+1\). Since \(t+2<2t+1\), no pair suffices.
Substitution gives the claimed difference.
\end{proof}

\section{Why eleven vertices are necessary}
The graph \(G_5\) has 11 vertices, 20 edges, and \(\gtot M=3\cdot7=21\),
whereas \(2|C|=22\). DeLaVi\~na's catalogue reports that all connected
graphs through order ten were previously checked. The following proof
establishes the same lower bound directly.

\begin{theorem}
Every counterexample to Conjecture 255 has at least 11 vertices.
\end{theorem}
\begin{proof}
Put \(s=\gtot(G)\), \(M=M(G)\), and \(\Delta=\max_v |N(v)|\).
In a counterexample \(sM<2|C|\le2n\). If \(s=2\), a total dominating
edge has neighborhood union \(V\), giving \(sM=2n\).

If \(s=3\), a minimum total dominating set induces a graph on three
vertices without isolates. Label them \(u,v,w\) with \(uv,vw\in E\).
The two edge neighborhoods cover \(V\). Each contains all of \(u,v,w\),
so \(n\le 2M-3\). It follows that \(3M<2n\le4M-6\), hence \(M\ge7\)
and \(2n>3M\ge21\).

It remains to exclude \(s\ge4\) and \(n\le10\) simultaneously.
These inequalities give \(4M\le sM<2n\le20\), so \(M\le4\).
A universal vertex and any one neighbor form a total dominating pair,
so no vertex is universal. Take a vertex \(v\) of degree \(\Delta\) and,
by connectivity, a vertex \(z\) at distance two along a path \(v,w,z\).
The union \(N(v)\cup N(w)\) contains \(N(v)\) as well as \(v\) and \(z\).
Thus \(M\ge\Delta+2\) and \(\Delta\le2\), which yields \(M\ge2\Delta\).
The \(s\) open neighborhoods of a total dominating set cover \(V\), so
\(s\Delta\ge n\). Now \(sM\ge2n\), a contradiction.
\end{proof}

\section{The sharp coefficient}
\begin{theorem}
Over all finite simple connected graphs of order at least two,
\[
 \inf_G\frac{\gtot(G)M(G)}{|C(G)|}=1.
\]
\end{theorem}
\begin{proof}
An edge incident with a vertex of maximum degree shows \(M\ge\Delta\).
The covering bound \(s\Delta\ge n\) therefore gives \(sM\ge n\ge|C|\).

For the reverse inequality fix \(k\ge2\), set \(q=k^2\), and start with
the \(q\)-by-\(q\) rook graph (cells sharing a row or a column are adjacent).
Partition the rows into \(k\) groups of \(k\) rows. For each group add a
hub joined to every cell in those rows, and add an apex \(x\) joined to
all hubs. These are all the edges of the resulting graph \(H_k\).

Two cells have distance at most two on the board. From a hub to a cell
outside its group, use any cell of the group in the target column.
Two hubs share \(x\), and \(x\) reaches every cell through its hub.
No vertex is universal, so \(|C(H_k)|=|V(H_k)|=k^4+k+1\).
The apex and all hubs form a total dominating set of size \(k+1\).
The degrees of a cell, a hub, and the apex are respectively
\(2k^2-1,\ k^3+1,\ k\). Hubs are pairwise nonadjacent. On every edge
one endpoint therefore has degree at most \(2k^2-1\), and the other
has degree at most \(k^3+1\). Bounding the union by the degree sum gives
\[
 1\le\frac{\gtot(H_k)M(H_k)}{|C(H_k)|}
 \le\frac{(k+1)(k^3+2k^2)}{k^4+k+1}\longrightarrow 1 .
\]
Thus every coefficient strictly larger than one fails on some \(H_k\).
\end{proof}

\paragraph{Computational checks.}
The accompanying Python script reconstructs the graphs from their edge
lists, computes eccentricities by breadth-first search, and counts the
minimum total dominating sets of \(G_5\) by testing vertex subsets in
increasing size. It finds five such triples. It checks \(G_t\) for
\(2\le t\le12\), bounds on \(H_k\) for \(k=2,3,4\), and all 27,475
connected labeled graphs of orders two through six. The proofs above cover
the full parameter ranges.

\paragraph{Source.}\raggedright
Ermelinda DeLaVi\~na, \emph{Written on the Wall II: Conjectures of
Graffiti.pc}, University of Houston--Downtown, entry 255 and catalogue
introduction. The catalogue credits the original bounded computation;
its November 21, 2026 note credits Tony Newton with a 14-vertex
counterexample to entry 256, whose numerator uses the neighborhood of
the minimum-degree vertices. Searches by conjecture number and invariants on
GitHub, arXiv, and Zenodo, did not identify an earlier resolution of entry
255. The original catalogue is available at\par\noindent
\url{https://www.uhd.edu/academics/sciences/mathematics-statistics/wowii/index.html}.
\end{document}
